% Archived from commit 819142e on 2026-09-12; not included in main_rewrite.tex.
\section{Equilibrium without an upper bound on AI efficiency}
\label{sec:rewrite-uncapped}

I now remove the upper bound on AI efficiency and compare the three
substitution regimes. The research law becomes
\begin{equation}
 \dot B=\chi(BM)^\eta.
 \label{eq:rewrite-uncapped-unit-law}
\end{equation}
All other primitives remain as in Section~\ref{sec:rewrite-model}.
In Definition~\ref{def:rewrite-equilibrium}, I set $\psi(B)=1$ and
$\psi'(B)=0$ and remove the upper restriction on $B$. Equilibrium still
requires finite allocations at every finite date, global agent optimality,
and both transversality conditions. Removing the upper bound does not by itself
imply that $B$ grows without bound.
%For any positive variable $V$, I write $g_V=\dot V/V$ for its growth rate.

The results differ in scope, depending on the elasticity of substitution. Under complementarity, I prove that long-run output per worker is still given by the rate of labor-augmenting technological change. At unit elasticity, there is the need of parameter restrictions for an equilibrium to exists. Such an equilibrium is analytically tractable. Under substitution, I distinguish an unbounded developer objective from a conditional finite-time singularity; either prevents the corresponding path from being an equilibrium.

\subsection{Complementarity}
\label{subsec:rewrite-uncapped-complements}

When $sigma<1$, effective labor remains necessary even if AI becomes arbitrarily efficient. Since $sigma<1$,  Equation~\eqref{eq:rewrite-ces} gives an upper bound for $Z$ and $Y$ which does not depend on $B$:
\begin{align}
 Z&\leq\omega_L^{\sigma/(\sigma-1)}AL,
 \label{eq:rewrite-uncapped-complements-z-bound}\\
 Y&\leq\omega_L^{\sigma(1-\alpha)/(\sigma-1)}
 K^\alpha(AL)^{1-\alpha}.
 \label{eq:rewrite-uncapped-complements-y-bound}
\end{align}

\Needspace{7\baselineskip}
\begin{proposition}[Production bounds under complementarity]
\label{prop:rewrite-uncapped-complements-bounds}
Let $0<\sigma<1$ and $\overline B=\infty$.
Along any feasible path satisfying the production equations~\eqref{eq:rewrite-output}
and~\eqref{eq:rewrite-ces}, labor-market clearing $L=N$ in
equation~\eqref{eq:rewrite-eq-labor}, the uncapped research law~\eqref{eq:rewrite-uncapped-unit-law},
and the resource constraint~\eqref{eq:rewrite-resource}, the following holds:
\begin{equation}
 \limsup_{t\to\infty}\frac1t
 \log\left(\frac{Y_t/N_t}{Y_0/N_0}\right)\leq\gamma.
 \label{eq:rewrite-uncapped-complements-growth-bound}
\end{equation}
%Neither $K$, $Y$, nor $B$ can become unbounded at a finite date.
\end{proposition}
\begin{proof}
See Appendix~\ref{app:rewrite-proofs},
\hyperref[proof:rewrite-uncapped-complements-bounds]{Proof of
Proposition~\ref*{prop:rewrite-uncapped-complements-bounds}}.
\end{proof}

Thus a permanent growth rate of output per person above $\gamma$ is infeasible, although transitional growth can exceed it. %These are feasibility bounds, not an equilibrium-existence theorem. They also do not require $\eta<\alpha$: $0<\eta<1$ suffices to bound AI efficiency on a finite horizon once research expenditure is bounded by available resources.
Complementarity suffices to bound long-run growth, despite no bound on AI efficiency.

%The monopoly matters for the limiting income shares. If its marginal cost $1/B$ tends to zero, the developer does not supply unlimited AI services. It approaches the quantity that maximizes revenue. The following conditional result states the implications without assuming that such an equilibrium has been constructed.

\Needspace{7\baselineskip}
\begin{proposition}[Long-run outcomes under complementarity without an AI efficiency upper bound]
\label{cond:rewrite-uncapped-complements-limits}
Suppose $0<\sigma<1$, $\overline B=\infty$, and an equilibrium satisfies
$B\to\infty$, while $K/(AL)$ and $C/(AL)$ converge to positive finite
constants. If $g_C$, $g_Y$, and $g_w$ have finite limits, then
\begin{align}
 g_{Y/N}&\to\gamma,\qquad g_w\to\gamma,\qquad r\to\rho+\gamma,
 \label{eq:rewrite-uncapped-complements-rates}\\
 s_X&\to\frac{1-\sigma}{1-\alpha\sigma},
 \label{eq:rewrite-uncapped-complements-sx}\\
 \frac{wL}{Y}&\to\frac{\sigma(1-\alpha)^2}{1-\alpha\sigma}>0,
 \qquad
 \frac{p_XX}{Y}\to\frac{(1-\alpha)(1-\sigma)}{1-\alpha\sigma}>0.
 \label{eq:rewrite-uncapped-complements-income}
\end{align}
\end{proposition}
\begin{proof}
See Appendix~\ref{app:rewrite-proofs},
\hyperref[proof:rewrite-uncapped-complements-limits]{Proof of
Proposition~\ref*{cond:rewrite-uncapped-complements-limits}}.
\end{proof}

%This proposition characterizes the long-run outcomes of an equilibrium satisfying the stated convergence conditions; it does not establish its existence.

%The limiting income shares follow from the static monopoly condition, equation~\eqref{eq:rewrite-monopoly-foc}. At positive limiting capital per unit of effective labor, $1/B\to0$ implies $e_X\to1$. Substituting into equation~\eqref{eq:rewrite-inverse-demand-elasticity} gives \eqref{eq:rewrite-uncapped-complements-sx}; equations~\eqref{eq:rewrite-wage} and~\eqref{eq:rewrite-ai-price} then give the income shares. The industry retains positive revenue even though its inference cost vanishes relative to output.

\subsection{Unit elasticity}
\label{subsec:rewrite-uncapped-unit-bgp}

At unit elasticity, define the output elasticities
\begin{equation}
 \beta=(1-\alpha)\omega_X,\qquad
 \lambda=(1-\alpha)\omega_L,
 \qquad\alpha+\beta+\lambda=1.
 \label{eq:rewrite-uncapped-unit-elasticities}
\end{equation}
Production is then $Y=K^\alpha(AL)^\lambda X^\beta$.
Consider a balanced-growth path (BGP) with positive constant $K/Y$, $C/Y$,
and $M/Y$. The AI-demand and
monopoly conditions, equations~\eqref{eq:rewrite-ai-price}
and~\eqref{eq:rewrite-monopoly-foc}, imply $U/Y=\beta^2$.
Using $X=BU$ in production, and $g_M=g_Y$ in the research law, gives
\begin{align}
 g_Y&=n+\gamma+\frac{\beta}{\lambda}g_B,
 \label{eq:rewrite-uncapped-unit-production-growth}\\
 (1-\eta)g_B&=\eta g_Y.
 \label{eq:rewrite-uncapped-unit-research-growth}
\end{align}
The first relation describes how AI improvement raises production growth;
the second describes how expanding research resources sustain AI improvement.
Solving these relations gives candidate growth rates. The next proposition
establishes that they belong to an equilibrium.

\Needspace{7\baselineskip}
\begin{proposition}[An uncapped balanced-growth equilibrium]
\label{prop:rewrite-uncapped-unit-bgp}
Suppose $\sigma=1$, $\overline B=\infty$, and
\begin{equation}
 0<\eta<\alpha,\qquad \eta\leq\tfrac12,\qquad
 \beta\leq\tfrac12,\qquad \eta+\omega_X<1.
 \label{eq:rewrite-uncapped-unit-sufficient}
\end{equation}
For every $A_0,N_0>0$, there exist $K_0,B_0>0$ and an equilibrium such that
\begin{align}
 g_{Y/N}=g_w& =\gamma+\frac{\eta\omega_X}{1-\eta-\omega_X}(n+\gamma),
 \label{eq:rewrite-uncapped-unit-wage-growth}\\
 r&=\rho+\gamma+\frac{\eta\omega_X}{1-\eta-\omega_X}(n+\gamma),
 \label{eq:rewrite-uncapped-unit-return}\\
 \frac{wL}{Y}&=\lambda=(1-\alpha)\omega_L>0,
 \label{eq:rewrite-uncapped-unit-labor-share}\\
 g_B&=\frac{\eta(1-\omega_X)}{1-\eta-\omega_X}(n+\gamma).
 \label{eq:rewrite-uncapped-unit-efficiency-growth}
\end{align}
%Both agents attain finite-valued global optima; both transversality conditions hold.
\end{proposition}
\begin{proof}
See Appendix~\ref{app:rewrite-proofs},
\hyperref[proof:rewrite-uncapped-unit-bgp]{Proof of
Proposition~\ref*{prop:rewrite-uncapped-unit-bgp}}.
\end{proof}

This is balanced growth, not a steady state in levels. The proposition
constructs the initial stocks together with the equilibrium; it does not
establish existence for arbitrary $K_0,B_0$. Appendix~\ref{app:rewrite-uncapped-unit}
derives all allocation ratios and states sufficient conditions for
equilibrium trajectories from nearby stocks that need not lie on the BGP.
The condition $\eta+\omega_X<1$ is necessary for this positive finite-rate
BGP. The bounds $\eta\leq1/2$ and $\beta\leq1/2$ are sufficient for the
global optimality proof, not necessary conditions for equilibrium existence.
The appendix explains the role of each restriction.

With any finite upper bound, the characterized unit-elastic equilibria have
$g_{Y/N},g_w\to\gamma$. Without an upper bound, continued AI improvement can
sustain faster growth of both quantities while labor retains a positive
income share. Thus the equivalence between a wage-growth premium and a
vanishing labor share across the regimes with a finite upper bound does not extend
to the uncapped economy. In the companion working paper
\emph{How Much Can Recursive AI Self-Improvement Raise Economic Growth?}
\citep{pardo2026companion}, I develop this unit-elastic benchmark as a framework
for calibration and sensitivity analysis of long-run growth.

\subsection{Substitution: unbounded returns and equilibrium existence}
\label{subsec:rewrite-uncapped}

Proposition~\ref{prop:rewrite-equilibrium-regimes} does not establish equilibrium existence without an upper bound for AI efficiency. I first fixed $\overline B$ and characterized an equilibrium for initial conditions near the steady state. Increasing $\overline B$ afterward differs from setting $\overline B=\infty$ at the outset. For $\sigma>1$, the limiting growth and interest rates of the AI-dominated equilibria become unbounded as $\overline B\to\infty$. This divergence is not a proof of a finite-time singularity or of equilibrium nonexistence in the uncapped economy.

Without an upper bound, recursive improvement can make the developer's objective unbounded: higher $B$ raises both operating profit and research productivity. If $\sigma>1$ and $\eta>\alpha$, there is no uncapped equilibrium with a finite-valued developer optimum, as established in Proposition~\ref{prop:rewrite-research-scale} in Appendix~\ref{app:rewrite-proofs}. Expanding early research makes discounted net profit arbitrarily large on a fixed finite horizon. The developer can then resume any proposed continuation of research spending with higher AI efficiency and no lower operating profit. Thus every candidate with a finite value admits a profitable deviation. This failure of optimality does not require a finite-time singularity: each deviating plan uses finite research expenditure and keeps $B$ finite on every finite horizon. %Financing cannot resolve an unbounded objective or create the real resources used in research: equation~\eqref{eq:rewrite-eq-kdot} requires research compute to compete with consumption, inference compute, and physical investment.

If $\eta<\alpha$ and $\sigma>1$, the same proposition gives a finite truncated value and shows that research costs eventually outgrow the benefits of increasing expenditure on a fixed horizon. This result takes the paths of capital, effective labor, and interest rates as given. It rules out that particular unbounded-payoff argument, without establishing an uncapped infinite-horizon equilibrium. Neither existence nor general nonexistence is established for this parameter range. On its own, this fixed-horizon result does not establish or rule out a finite-time singularity, which would violate the requirement of finite, feasible allocations at every date in Definition~\ref{def:rewrite-equilibrium}.

The equilibrium in Section~\ref{subsec:rewrite-uncapped-unit-bgp} shows that
unbounded AI efficiency is compatible with equilibrium at unit elasticity.
It does not resolve existence for $\sigma>1$.

At $\eta=\alpha$, the finite-horizon benefit and cost terms have the same
order, so their coefficients determine whether that particular
unbounded-payoff argument applies. Proposition~\ref{prop:rewrite-research-scale}
does not settle equilibrium existence at this boundary.

The literature suggests why $\eta<\alpha$ need not prevent an explosion.
\citet{nordhaus2021} emphasizes the role of substitution in singular growth,
and \citet{davidsonetal2026} show how technological and economic feedback
can jointly overcome diminishing returns in research. Their capital-accumulation
analysis holds saving and factor-allocation shares fixed. I therefore cannot
apply its explosion theorem directly to this model, where those allocations
are endogenous. The following result instead uses this model's production,
monopoly, research, and household Euler equations.

\Needspace{7\baselineskip}
\begin{proposition}[Conditional finite-time explosion under substitution]
\label{prop:rewrite-uncapped-substitutes-explosion}
Let $\sigma>1$, $\overline B=\infty$, $0<\alpha<1$, and $0<\eta<1$.
Consider a differentiable solution of the intratemporal and dynamic
equations~\eqref{eq:rewrite-equilibrium-static} and~\eqref{eq:rewrite-equilibrium-dynamics}
in Definition~\ref{def:rewrite-equilibrium}, with $K,B,C,M>0$ and research
law~\eqref{eq:rewrite-uncapped-unit-law}, on a maximal interval $[0,T)$,
where $T$ may initially be infinite.
Suppose that, as $t\uparrow T$, $B\to\infty$, $s_X\to1$, and $C/Y$ and
$M/Y$ converge to positive finite constants, with
\begin{equation}
 \frac K Y\frac{\mathrm d}{\mathrm dt}\log\left(\frac C Y\right)\to0,
 \qquad
 \frac K Y\frac{\mathrm d}{\mathrm dt}\log\left(\frac M Y\right)\to0.
 \label{eq:rewrite-uncapped-explosion-regularity}
\end{equation}
Then
\begin{equation}
 \frac{\mathrm d(Y/K)/\mathrm dt}{(Y/K)^2}
 \longrightarrow\frac{\eta\alpha(1-\alpha)}{1-\alpha\eta}>0.
 \label{eq:rewrite-uncapped-explosion-coefficient}
\end{equation}
Consequently, $T<\infty$ and $B,K,Y,C,r\to\infty$ as $t\uparrow T$.
No equilibrium in Definition~\ref{def:rewrite-equilibrium} can have these
asymptotics, whether $\eta<\alpha$, $\eta=\alpha$, or $\eta>\alpha$.
\end{proposition}
\begin{proof}
See Appendix~\ref{app:rewrite-proofs},
\hyperref[proof:rewrite-uncapped-substitutes-explosion]{Proof of
Proposition~\ref*{prop:rewrite-uncapped-substitutes-explosion}}.
\end{proof}

Condition~\eqref{eq:rewrite-uncapped-explosion-regularity} requires the
proportional changes in consumption and research shares to become negligible
relative to $Y/K$. Convergence of the shares alone would not imply this
condition. Neither constant shares throughout the transition nor convergence
of $g_B/(Y/K)$ is assumed; the proof derives the latter limit.

The mechanism is cumulative. When AI dominates the composite input, better AI
raises output per unit of capital. The higher return raises consumption
growth through the Euler equation. If consumption and research retain
positive, sufficiently stable shares of output, research spending expands
with output and sustains further AI improvement. This feedback makes $Y/K$
grow in proportion to its own square, which implies a finite-time singularity.
The restriction $\eta<\alpha$ limits the profitability of increasing research
at given prices; it does not eliminate this joint feedback.

Proposition~\ref{prop:rewrite-uncapped-substitutes-explosion} is conditional:
it does not construct a path satisfying its hypotheses or show that every
initial condition leads to them. A solution of the dated equations is not
necessarily an equilibrium, because global optimality and the transversality
conditions must also hold. Here finite-time divergence already precludes an
infinite-horizon equilibrium. For $\eta>\alpha$, the separate profitable-deviation
argument in Proposition~\ref{prop:rewrite-research-scale} rules out a
finite-valued developer optimum without these convergence assumptions.
For $\eta\leq\alpha$, existence or nonexistence outside the asymptotic class
excluded above remains open.

